% !TEX program = xelatex
\documentclass{英语试卷}
\begin{document}
\examtitle{2024 学年第一学期浙江省县域教研联盟高三年级模拟考试}{英语}
\par 英语学科 试题卷\par
\par 命题：临海市教研中心 审题：海宁市教师进修学校 柯桥区教师发展中心\par
\begin{notice}
\par 1. 答卷前，考生务必将自己的姓名、准考证号填写在答题卡上，\par
\par 2. 回答选择题时，选出每小题答案后，用铅笔把答题卡上对应题目的答案标号涂黑。如需改动，用橡皮擦干净后，再选涂其他答案标号。回答非选择题时，将答案写在答题卡上，写在本试卷上无效，\par
\par 3. 考试结束后，将本试卷和答题卡一并交回。\par
\end{notice}
\parttitle{第一部分 听力（共两节，满分 30 分）}
\par 做题时，先将答案标在试卷上，录音内容结束后，你将有两分钟的时间将试卷上的答案转涂到答题卡上。\par
\sectitle{第一节（共5小题：每小题1.5分，满分7.5分）}
\par 听下面 5 段对话。每段对话后有一个小题，从题中所给的 A、B、C 三个选项中选出最佳选项。听完每段对话后，你都有 10 秒钟的时间来回答有关小题和阅读下一小题。每段对话仅读一遍。\par
\question{1. What will the man do tonight?}
\begin{choices}(3)
  \item Offer a recipe.
  \item Make some cookies.
  \item Go to a bakery.
\end{choices}
\question{2. What are the speakers talking about?}
\begin{choices}(3)
  \item A vacation plan.
  \item Weather forecast.
  \item After-class activities.
\end{choices}
\question{3. What does the woman want to watch?}
\begin{choices}(3)
  \item Evening news.
  \item A sports game.
  \item A quiz show
\end{choices}
\question{4. How does the woman feel about the store?}
\begin{choices}(3)
  \item The prices are low.
  \item It has limited choices.
  \item The service needs improving.
\end{choices}
\question{5. What is probably the relationship between the speakers?}
\begin{choices}(3)
  \item Teacher and student.
  \item Schoolmates.
  \item Shop assistant and owner.
\end{choices}
\sectitle{第二节（共15小题：每小题1.5分，满分22.5分）}
\par 听下面 5 段对话或独白。每段对话或独白后有几个小题，从题中所给的 A、B、C 三个选项中选出最佳选项。听每段对话或独白前，你将有时间阅读各个小题，每小题 5 秒钟；听完后，各小题将给出 5 秒钟的作答时间。每段对话或独白读两遍。\par
\question{6. What does the woman say about TV commercials?}
\begin{choices}(3)
  \item They ran for 5 months.
  \item They are successful.
  \item They cost a lot.
\end{choices}
\question{7. What does the man suggest?}
\begin{choices}(3)
  \item Having a meeting.
  \item Talking with the president.
  \item Learning from the competitors.
\end{choices}
\par 听第 7 段材料，回答第 8、9 题，\par
\question{8. What surprised the man about the room?}
\begin{choices}(3)
  \item The view.
  \item The cost.
  \item The service.
\end{choices}
\question{9. What will be extra charged?}
\begin{choices}(3)
  \item Room service.
  \item Use of the hotel's safe.
  \item A full continental buffet.
\end{choices}
\par 听第 8 段材料，回答第 10 至 13 题。\par
\question{10. What tour does the man decide to take today?}
\begin{choices}(3)
  \item A three-hour tour.
  \item A half-day tour.
  \item A one-day tour.
\end{choices}
\question{11. What time is it now?}
\begin{choices}(3)
  \item 10:00 a. m
  \item 10:30 a. m.
  \item 11:00 a. m.
\end{choices}
\question{12. Which place will the man visit tomorrow?}
\begin{choices}(3)
  \item Buckingham Palace.
  \item Windsor Castle.
  \item Westminster Abbey.
\end{choices}
\question{13. How much does the man pay in all?}
\begin{choices}(3)
  \item £15.
  \item £90.
  \item £105.
\end{choices}
\par 听第 9 段材料，回答第 14 至 17 题\par
\question{14. Where did the speakers see Brigette Clark, the famous movie star?}
\begin{choices}(3)
  \item In the Sandwich Place.
  \item In the American style diner.
  \item In the new Italian restaurant.
\end{choices}
\question{15. What can be learned about the American style diner?}
\begin{choices}(3)
  \item They serve quickly.
  \item They make large hamburgers.
  \item They sell good chips.
\end{choices}
\question{16. How did the man know the Italian restaurant was great?}
\begin{choices}(3)
  \item He tried its food himself.
  \item He read reviews online.
  \item He heard it from his wife.
\end{choices}
\question{17. What did the woman think of Jeremy's lunch?}
\begin{choices}(3)
  \item It's delicious.
  \item It's nutritious.
  \item It's awful.
\end{choices}
\par 听第 10 段材料，回答第 18 至 20 题，\par
\question{18. What did people worry about concerning the post office?}
\begin{choices}(3)
  \item It would be closed down.
  \item It would be expensive to run.
  \item It would be moved out of Scotland.
\end{choices}
\question{19. Why did the Fords want to come back home?}
\begin{choices}(3)
  \item To start a company.
  \item To look after their mothers.
  \item To do business with a company
\end{choices}
\question{20. What do the Fords expect?}
\begin{choices}(3)
  \item To receive more visitors.
  \item To build a new community.
  \item To cut back postal services.
\end{choices}
\parttitle{第二部分 阅读理解（共两节，满分 50 分）}
\sectitle{第一节（共15小题：每小题2.5分，满分37.5分）}
\par 阅读下列短文，从每题所给的 A、B、C 和 D 四个选项中，选出最佳答案，\par
\passagelabel{A}
\par Brogdale Collections is the home of the National Fruit Collection and works to provide access and education about the National Fruit Collection. It offers a variety of activities to suit all ages and interests. For upcoming events, look here.\par
\par Heritage Orchard Days\par
\par Wed, October 9, 2024-Sun, October 20, 2024(1:00 pm-3:00 pm)\par
\par Come and join us for a fun family day out at Brogdale Farm. Each weekend, visitors will delight in the wonderful orchards(果园) bursting with fruit, join a guided tour of the National Collection and taste hundreds of varieties straight from the trees. Children will join in with some apple campfire cooking and traditional apple pressing. Included in the ticket price, visitors will also harvest their own apples and pears to take home.\par
\par Ticket price: Adult(16yrs+)£10.00 Child(3-15yrs, under 3's free)£5.00\par
\par Kids Juicing Workshops Mon, August 19, 2024-Mon, October 21, 2024(10:00 am-3:00 pm)\par
\par Everyone loves apple juice, but have you ever thought about where it comes from and how it is made? Here we are giving children a unique opportunity to explore where our apples are growing and find out what happens to them. Come along to Brogdale where you will see more than 3000 varieties of fruit trees including apples, cherries, pears and plums. Then, there is an opportunity to pick some apples and have a go at making freshly pressed juice!\par
\par Ticket price: £15.00 per child plus accompanying adult Additional child £10.00 Additional adult£15.00\par
\question{21. What is the main purpose of Brogdale Collections?}
\begin{choices}(2)
  \item Offering traditional gardening courses.
  \item Collecting various fruits from all states.
  \item Providing fruit-related education with fun.
  \item Making tailored package tours for families.
\end{choices}
\question{22. What do the two events have in common?}
\begin{choices}(2)
  \item They offer guided tours.
  \item They end on the same day.
  \item They involve apple pressing.
  \item They allow making campfire.
\end{choices}
\question{23. How much do a couple with a kid pay to join Kids Juicing Workshop?}
\begin{choices}(2)
  \item £25.
  \item £30.
  \item £40.
  \item £45.
\end{choices}
\passagelabel{B}
\par Street artist Banksy, whose real identity remains unknown, has grabbed public attention with animals painted at different spots around London this week. On Aug. 5, a water bird was drawn above a fish-and-chips shop, as though wolfing down fish on the store's signs. Then the following day, monkeys were seen painted swinging across a bridge near an East London station. "It's an eye trick, and that's what draws you in," says Tordoff, a specialist in such art.\par
\par There is no clear explanation for this animal-focused art series. Some have thought of it as a warning that eventually humanity would lose control of the earth and cities like London could return to nature. Meanwhile, a BBC presenter has claimed that Banksy cares deeply about the environment, animal welfare and ecological issues. "He puts situations together which are at first, playful and humorous," says Acoris Andipa, who runs a gallery and sells Banksys worldwide. "People laugh, and then suddenly, step back and just say, 'Ouch, that's got a sting to it.' That's something people really enjoy."\par
\par Banksy's work challenges the art market — often adding to its charm. Andipa has sold hundreds of Banksy's paintings and thousands of prints since the mid-2000s. He mentioned that besides art thefts, some pieces have been ruined by other graffiti artists since their appearance. But with his knowledge of the commercial market for Banksy's work, he says the recent London series is unlikely to influence prices for his work significantly.\par
\par Almost everyone in the art world has a theory about Banksy's real identity. "He's hiding behind this mask: If his identity was revealed, immediately you would associate his painting with his position. It would be easy to say, I don't like it." Tordoff says. Banksy published his works on Instagram, allowing him to reach a large global audience — with more than 13 million followers.\par
\question{24. According to Tordoff, what is the magic of Banksy's work?}
\begin{choices}(2)
  \item Its mystery.
  \item Its visual effect.
  \item Its location.
  \item Its animal theme
\end{choices}
\question{25. What does paragraph 2 focus on?}
\begin{choices}(2)
  \item Influence of Banksy's work.
  \item Views about Banksy's work.
  \item Concerns over Banksy's work.
  \item Suggestions on Banksy's work.
\end{choices}
\question{26. What does Andipa think might happen to Banksy's work?}
\begin{choices}(2)
  \item Its price will stay relatively stable.
  \item It will be attacked by other artists.
  \item Its commercial market will expand.
  \item It will continue to focus on animals.
\end{choices}
\question{27. What is Tordoff's attitude to Banksy hiding his identity?}
\begin{choices}(2)
  \item Critical.
  \item Ambiguous.
  \item Understanding.
  \item Doubtful.
\end{choices}
\passagelabel{C}
\par Halloween plays on our fears and our fantasies. We craft haunted houses and scary decorations to arouse particular emotions. We choose our costumes to reflect the kind of people we are or want to be — edgy, funny, clever. For children, Halloween is an experiment in delayed satisfaction and negotiation — which candies to eat now, which to trade, which to save. It's no surprise, then, that Halloween might show interesting features of human psychology.\par
\par In fact, there's a long tradition of using Halloween to shed light on the human mind and behavior. In a study published recently, researchers observed over 1,000 trick-or-treating children as they visited houses in Seattle on the evening of Halloween. The researchers were interested in understanding the conditions that lead to "unrestricted" behavior: stealing Halloween candy or money. One of the variables they controlled was whether the adult who greeted the children at the entrance asked for each child's name and address, thereby treating each child as an identifiable individual, or instead let each child remain a stranger. Either way, the adult then instructed each child to take one candy from the table while the adult went away to "work in another room."\par
\par The children's behavior was recorded by an observer behind a peephole. For each child, the observer recorded how many candies were taken, and whether the child took any money from a bowl of coins next to the candy. And they did take candy and money: about 30 percent of children took extra candy, money or both.\par
\par The researchers identified several factors that influenced the probability of such behavior. Thefts were more likely for children whose names remained unknown, who were in groups rather than alone, and who were not accompanied by an adult. There was also an important influence of peer behavior: kids in groups were more likely to steal if the first child in their group did so.\par
\question{28. Which of the following might children care more on Halloween?}
\begin{choices}(2)
  \item Horrible houses.
  \item Funny costumes.
  \item Scary decorations.
  \item Trick-or-treat candies.
\end{choices}
\question{29. What was the purpose of the Halloween research?}
\begin{choices}(2)
  \item To renew the knowledge of Halloween.
  \item To give a new definition to "trick or treat".
  \item To uncover a Halloween-related psychology.
  \item To expose children's bad behavior on Halloween.
\end{choices}
\question{30. How was the research mainly conducted?}
\begin{choices}(2)
  \item By observing subjects.
  \item By analyzing data.
  \item By recording children's identities.
  \item By controlling children's behavior.
\end{choices}
\question{31. Which group of children is less likely to steal candies?}
\begin{choices}(2)
  \item Kids who remain nameless.
  \item Kids coming with their parents
  \item Kids accompanied by their friends.
  \item Kids whose peer members steal candies.
\end{choices}
\passagelabel{D}
\par It is a common view that "Necessity is the mother of invention." That is, inventions supposedly arise when a society has an unfulfilled need. Quite a few inventions do fit in this category. For example, in 1794 Eli Whitney invented cotton gins to replace laborious hand cleaning of cotton grown in the U. S. South.\par
\par Such familiar examples fool us into assuming that other major inventions were also responses to perceived needs. In fact, most inventions were developed by people driven by curiosity or by a love of tinkering (捣鼓). Once a device had been invented, the inventor then had to find an application for it. Only after it had been in use for a considerable time did consumers come to feel that they "needed" it. Thus, invention is often the mother of necessity, rather than vice versa.\par
\par A good example is the history of the motor vehicle which was not invented in response to any demand. When Nikolaus Otto built his first gas engine in 1866, it was weak, heavy, and seven feet tall. Thirty years later, he built the first truck. But it was a time when horse wagons and steam-powered railroads dominated transportation. Public contentment with these two means remained high until World War I when the armies concluded that they really did need trucks, which eventually made those vehicles a substitute for horse-drawn wagons in industrialized countries.\par
\par Inventors often have to persist at their tinkering for a long time in the absence of public demand, because early models perform too poorly to be useful. For instance, the first cameras, typewriters, and television sets were as awful as Otto's seven-foot-tall gas engine. That makes it difficult for an inventor to foresee whether his or her awful prototype might eventually find a use and thus invest more time and expense to develop it. Even inventions that meet the need for which they were initially designed may later prove more valuable at meeting unforeseen needs. While James Watt designed his steam engine to pump water from mines, it soon was supplying power to cotton mills, then(with much greater profit) propelling trains and boats.\par
\question{32. Why does the writer mention Eli Whitney in the first paragraph?}
\begin{choices}(2)
  \item To clarify a doubt.
  \item To illustrate a view.
  \item To introduce the topic
  \item To evaluate a statement.
\end{choices}
\question{33. What contributed to the popularity of motor vehicles?}
\begin{choices}(2)
  \item Military demands.
  \item Powerful engines.
  \item Creative designs.
  \item Public satisfaction.
\end{choices}
\question{34. What does the underlined word "That" refer to in the last paragraph?}
\begin{choices}(2)
  \item A lack of money and time.
  \item Absence of public demand.
  \item Bad performance of early models.
  \item Awful looks of previous inventions.
\end{choices}
\question{35. Which of the following is the best title for the text?}
\begin{choices}(2)
  \item Social Need: the Inner Drive for Invention
  \item Great Invention: the Force of Society's Progress
  \item General Application: the Final Purpose of Invention
  \item Beyond Necessity: the Curious Pursuit of Invention
\end{choices}
\sectitle{第二节（共5小题；每小题2.5分，满分12.5分）}
\par 阅读下面短文，从短文后的选项中选出能填入空白处的最佳选项。选项中有两项为多余选项。\par
\par It's your turn to talk, your opportunity to dive deep into a favorite anecdote. People around you are listening, hanging on your every word — until someone interrupt. There is little for you to do apart from swallow the bitter taste in your mouth, smile, nod and return to listening mode. \blanknum{36} Perhaps not.\par
\par First thing first. Although it might be easy to imagine you are not interesting or important enough to speak, that isn't the case. In my opinion, one doesn't have to get to the point or avoid boring topics all the time to be a good speaker. \blanknum{37}\par
\par So when people seem to brush aside your story next time, how to respond appropriately? \blanknum{38} Everyone is programmed slightly differently. Perhaps if it's a first-time incident, just take a deep breath and let it go. That's not to say you can't look out for a chance to return to your anecdote, just that now's time to press the pause button.\par
\par However, if you don't want the interrupter to speak over you, there are gracious ways to reclaim your voice \blanknum{39} Insist you would love to hear their thoughts when you've finished your story. Then, in theory, everyone wins.\par
\par \blanknum{40} In group environments, that meaningful grin from across the table, the person with the "go-on" eyes, is waiting to hear your tale's ending. This might be because they're interested, or they recognize you've been silenced and want to show solidarity. Whatever it is, if you can, keep going.\par
\par Pretty much everyone will experience conversational clashes at some point. Remember, the art of conversation is the art of hearing as well as being heard.\par
\begin{choices}(1)
  \item Does it have to be this way?
  \item Can't they just wait until you finish?
  \item One reason might be understanding.
  \item Understanding is almost certainly the answer.
  \item Everyone has the right to make their voice heard.
  \item Politely but firmly express that you appreciate their input.
  \item Don't give up before checking if you still have someone's attention.
\end{choices}
\parttitle{第三部分 语言知识运用（共两节，满分 30 分）}
\sectitle{第一节（共15小题；每小题1分，满分15分）}
\par 阅读下面短文，从短文后各题所给的 A、B、C 和 D 四个选项中，选出可以填入空白处的最佳选项。\par
\par I'm a photographer. Sometimes my travels take me to places where people have much less than I do. Sometimes I feel a wave of \blanknum{41}, knowing that the photo I've just taken is now in my possession, but the person in the photo is left with only a \blanknum{42} of the image on the back of my camera.\par
\par A decade ago, while planning a trip to Thailand, I came up with a solution to \blanknum{43} ny guilt: a portable photo printer. Equipped with one, I was \blanknum{44} to hand out snapshots wherever I went.\par
\par Near Mai Sot, I walked into a city dump where many poor people \blanknum{45} among the mountains of trash. A family squatting in their open hut watched me \blanknum{46} as if I came to spy on them. I asked the father if I could take their photo. He shrugged, \blanknum{47} why I wanted to do so. My printer began to work. The moment the first print appeared, they approached and I was \blanknum{48}.\par
\par Soon I began making friends. Children tugged at my shirt wanting to be the next \blanknum{49}. Old ladies tried to \blanknum{50} the children so they could cut in line. I couldn't make the photos fast enough, so I showed a teenage boy how to \blanknum{51} machine.\par
\par As their enthusiasm grew, I noticed that many held their small photos \blanknum{52}, as if they were gazing at a(n) \blanknum{53} they'd never seen before.\par
\par After many such journeys, I was finally free of that bad feeling and grew \blanknum{54}. I'm a photographer, one who loves taking portraits and \blanknum{55} them, and spreading happiness in the most unexpected places.\par
\cz{41}{guilt}{panic}{pride}{satisfaction}
\cz{42}{description}{copy}{glimpse}{reflection}
\cz{43}{hide}{prove}{express}{ease}
\cz{44}{lucky}{ready}{hesitant}{honored}
\cz{45}{marched}{wandered}{searched}{sheltered}
\cz{46}{eagerly}{silently}{suspiciously}{bravely}
\cz{47}{expecting}{knowing}{wondering}{noticing}
\cz{48}{encouraged}{surrounded}{forgiven}{recognized}
\cz{49}{role}{model}{fan}{member}
\cz{50}{take over}{run after}{calm down}{drive away}
\cz{51}{operate}{activate}{polish}{protect}
\cz{52}{in disgust}{with patience}{in shame}{with awe}
\cz{53}{image}{shape}{pattern}{scene}
\begin{choices}(1)
  \item concerned
  \item sympathetic
  \item relieved
  \item grateful 55.
  \item revealing
  \item distributing
  \item exhibiting
  \item cherishing
\end{choices}
\sectitle{第二节（共10小题；每小题1.5分，满分15分）}
\par 阅读下面短文，在空白处填入 1 个适当的单词或括号内单词的正确形式。\par
\par When international travelers began arriving in China in early 2023 following a three-year hiatus (停滞), many may have found that routine daily tasks \blanknum{56} (pose) challenges. More at ease with cash and credit cards and unfamiliar with Chinese apps, they experienced difficulties \blanknum{57} (perform) tasks such as paying with QR codes and booking tickets online.\par
\par But China is making efforts to address the \blanknum{58} (convenience) many international visitors may face during their stay. New initiatives are making travel around China \blanknum{59} (easy) in comparison to previous period. Alipay and Wexin Pay, China's digital payment services, now allow users to link international credit cards and make mobile payments in a way similar \blanknum{60} locals. Meanwhile, China is varying its means of payment to provide services \blanknum{61} (design) for foreign visitors. An increasing number of businesses upgrade their payment systems \blanknum{62} begin accepting international credit cards.\par
\par While the major ride-hailing app Didi Chuxing boasts \blanknum{63} user-friendly English interface(界面), museums and other attractions have been called on to develop websites \blanknum{64} provide multilingual booking services and maintain in-person ticketing services.\par
\par Since the introduction of a trial visa-free policy for visitors from certain countries, more countries \blanknum{65} (add) to the list. With continued efforts from the government as well as business and cultural sectors, travel experience in China will only become more enjoyable.\par
\parttitle{第四部分 写作（共两节，满分 40 分）}
\sectitle{第一节（满分15分）}
\par 假定你是李华，校英文报"Big Debate"栏目正面向全校学生征集关于未来职业的辩论话题。请给栏目编辑写一封邮件，内容包括：\par
\par (1) 推荐辩题;\par
\par (2) 说明理由。 注意：\par
\par (1) 写作词数应为 80 左右;\par
\par (2) 请按如下格式在答题纸的相应位置作答，\par
\par Dear Chris,\par
\par Yours, Li Hua\par
\sectitle{第二节（满分 25 分）}
\par 阅读下面材料，根据其内容和所给段落开头语续写两段，使之构成一篇完整的短文，\par
\par I hopped off the school bus, hitting the ground at a dead run, dropped my book bag at the house door and ran down to the barnyard. I couldn't wait to touch the little horse and continue the training. He was coming along so well!\par
\par My mother bought the horse as my 14th birthday gift. He was of unknown breeding, and a non-extraordinary brown color. The first time I caught his big, dark eyes with mine, I knew he was special. My mother was happy to see the two of us quickly bonded. As we progressed, his potential consumed me. He wasn't my first horse, but the first one I took on as a project and my first time to train a horse under saddle(马鞍).\par
\par However, my grin quickly faded into a confused frown as the little horse was no where to be seen. Panic surged through me. Where was my boy? Had something terrible happened? On the verge of tears, I dashed back to the kitchen for my mother, desperate for an answer.\par
\par "Mom! Where's ToiBoi? I can't find him anywhere!" My mother looked up from her chores as I burst through the door. "Oh, sweetie, come here," she said, reaching out to pull me to her side.\par
\par "I have some news," she began slowly. "You know how much we love and believe in you, and how proud we are of how far you and ToiBoi have come. But today something unexpected happened."\par
\par I stared at her, my mind racing with possibilities. "What happened, Mom? Is he hurt?"\par
\par "No, no," she assured me quickly. "He is not hurt. But there was an opportunity. A trainer came by this afternoon looking for promising young horses. He took ToiBoi to his training ground for test. He was very impressed and promised to give ToiBoi professional training. It's your horse. So I need you to make a choice for ToiBoi today."\par
\par (1) 续写词数应为 150 左右;\par
\par (2) 请按如下格式在答题卡的相应位置作答。\par
\par\vspace{0.4em}
\begin{center}\begin{tabular}{{|l|}}
\hline
Paragraph 1: "Professional training?" I echoed, trying to understand what she was saying. \\
 \\
 \\
Paragraph 2: Decision made, it was time to go and see ToiBoi in the training ground. \\
 \\
 \\
 \\
\hline
\end{tabular}\end{center}
\vspace{0.4em}
\end{document}
